\section{Related Work}
\label{sec:related}

\subsection{Single-image super-resolution}

Classical SR relied on interpolation and handcrafted priors. Deep CNN models
replaced these with learned mappings from low- to high-resolution
images~\cite{dong2014srcnn,lim2017edsr,zhang2018rcan}. Perceptual and
adversarial methods~\cite{ledig2017srgan,wang2021realesrgan} further increase
texture realism, often at the cost of higher entropy. Transformer-based
restorers such as SwinIR and Swin2SR~\cite{liang2021swinir,conde2022swin2sr}
currently provide strong accuracy--efficiency trade-offs; compact CNNs such as
NinaSR~\cite{gouvine2021torchsr} remain attractive when adaptation cost matters.
A common side effect across these families is that synthesised high-frequency
detail improves visual quality but can make the SR output harder to compress.

\subsection{Learned image compression}

Neural image compression (NIC) replaces hand-engineered transforms with
end-to-end trained nonlinear analysis/synthesis networks and learned entropy
models~\cite{balle2016end,balle2018variational,cheng2020learned}. Libraries
such as CompressAI~\cite{begaint2020compressai} make pretrained codecs easy to
reuse as differentiable bitrate interfaces. Most NIC work trains the codec
itself for rate--distortion performance. In contrast, we \emph{freeze} a
pretrained codec and use it only to measure and backpropagate bitrate while
adapting the upstream SR model.

\subsection{Compression-aware restoration and SR--compression pipelines}

The interaction between restoration and compression has been studied from
several angles. One line jointly trains restoration and compression networks
under a combined rate--distortion objective, so that both components co-adapt.
Another adds rate penalties to restoration losses while still optimising against
ground-truth high-resolution images. Related observations on generative and
perceptual SR note that richer textures raise bitrate under downstream
codecs~\cite{ledig2017srgan,blau2018perception}, but do not by themselves
provide a deployment-time adaptation procedure for an already pretrained
backbone.

These settings typically assume that both the restoration model and the codec
can be (re)trained, and that the optimisation target is ground-truth HR. In
many practical pipelines the codec is fixed by infrastructure constraints, and a
high-performing SR model is already deployed. Retraining the full system may be
undesirable or infeasible.

\subsection{Per-image and parameter-efficient adaptation}

Instance-specific optimisation and parameter-efficient fine-tuning
(e.g., LoRA~\cite{hu2021lora}) reduce the cost of adapting large networks to a
single input or a small parameter subset. We leverage this idea for
compression-aware SR: full fine-tuning of compact CNNs, full or LoRA-based
updates of Swin2SR, performed per image under a fixed codec.

\subsection{Positioning}

Table~\ref{tab:research_gap} summarises how the present setting differs from
common alternatives. Our problem is characterised by three choices:
(i)~a \emph{fixed} neural codec; (ii)~adaptation of a \emph{pretrained} SR
backbone rather than training from scratch; and (iii)~\emph{per-image}
optimisation whose fidelity reference is the baseline SR output, not the
ground-truth HR image. Constrained alternating optimisation via an
ADMM-inspired split~\cite{boyd2011admm} is used to keep bitrate minimisation
and fidelity preservation coupled and inspectable.

\begin{table}[t]
\centering
\caption{Comparison of optimisation settings in SR and image compression.}
\label{tab:research_gap}
\small
\begin{tabular}{@{}lcccc@{}}
\toprule
\textbf{Method} & \textbf{SR} & \textbf{Codec} & \textbf{GT} & \textbf{Img.-spec.} \\
\midrule
Standard SR              & Train  & Fixed  & Yes & No \\
Learned compression      & Fixed  & Train  & Yes & No \\
Joint SR+compression     & Train  & Train  & Yes & No \\
Compression-aware SR     & Train  & Fixed  & Yes & No \\
\textbf{Proposed}        & Adapt  & Fixed  & No  & Yes \\
\bottomrule
\end{tabular}
\end{table}
